\documentclass[11pt,a4paper]{article}

\usepackage[utf8]{inputenc}
\usepackage[T1]{fontenc}
\usepackage{babel}
\babelprovide[main, import]{french}
\usepackage[a4paper, margin=2.5cm]{geometry}
\usepackage{hyperref}
\usepackage{booktabs}
\usepackage{array}
\usepackage{listings}
\usepackage{xcolor}
\usepackage{amsmath}
\usepackage{amssymb}
\usepackage{textcomp}

\hypersetup{
    colorlinks=true,
    linkcolor=blue,
    urlcolor=blue,
    citecolor=blue
}

% Style pour coder l'assembleur RISC-V
\lstdefinestyle{riscv}{
    basicstyle=\ttfamily\small,
    keywordstyle=\color{blue}\bfseries,
    commentstyle=\color{gray}\itshape,
    numbers=left,
    numberstyle=\tiny\color{gray},
    breaklines=true,
    showstringspaces=false,
    tabsize=4,
    frame=single,
    backgroundcolor=\color{gray!5},
    morekeywords={addi,add,sub,mul,div,lw,sw,li,la,mv,j,jal,jalr,ret,bge,bgeu,blt,bltu,ble,beq,bne,bltz,beqz,blez,bgtz,slli,srli,srai,ecall,nop}
}

\lstdefinestyle{bash}{
    basicstyle=\ttfamily\footnotesize,
    breaklines=true,
    frame=single,
    backgroundcolor=\color{gray!5},
}

\title{Comparaison d'algorithmes de tri en RISC-V}
\author{Hubert Lestrat \\ \small Polytech-Sorbonne --- MAIN3 Architecture des ordinateurs}
\date{\today}

\begin{document}

\maketitle
\tableofcontents
\vspace{1em}

% =================================================================
\section{Références}
% =================================================================

\begin{itemize}
    \item \url{https://www.laurentbloch.net/MySpip3/Mon-premier-programme-en-assembleur-RISC-V}
    \item \url{https://www.laurentbloch.net/MySpip3/Nouvel-episode-de-programmation-en-assembleur-RISC-V}
    \item \url{https://www.laurentbloch.net/MySpip3/Assembleur-RISC-V-maintenant-les-tris}
    \item \url{https://fr.wikipedia.org/wiki/Tri_\%C3\%A0_bulles}
\end{itemize}

% =================================================================
\section{Astuces}
% =================================================================

\texttt{rars test.asm} $\Rightarrow$ compile et lance le programme en console.

Attention, lorsqu'on « include » un fichier, les étiquettes doivent rester uniques. Le mieux est de les préfixer avec le nom de la fonction.

Générer un tableau aléatoire de 20 entrées (pour \texttt{tab.asm}) :

\begin{lstlisting}[style=bash]
python3 -c "import random; random.seed(); \
  print(', '.join(str(random.randint(-99, 99)) for _ in range(20)))"
\end{lstlisting}

% =================================================================
\section{Méthode de mesures}
% =================================================================

\subsection{Dans RARS}

\begin{itemize}
    \item Ouvrir le fichier \texttt{.asm} à tester.
    \item \texttt{Tools} $\rightarrow$ \texttt{Instruction Counter}, puis « Connect to program ».
    \item \texttt{Settings} $\rightarrow$ \texttt{Hexadecimal Values}.
    \item \texttt{Settings} $\rightarrow$ \texttt{Hexadecimal Addresses}.
\end{itemize}

\begin{enumerate}
    \item \textbf{Assemble} (F3).
    \item Clic sur \textbf{Reset} dans la fenêtre Instruction Counter (important : vider entre deux tests).
    \item \textbf{Run} (F5) jusqu'à la fin.
    \item Relever les chiffres.
\end{enumerate}

\paragraph{Exécution.} Dans \texttt{Instruction Counter} :
\begin{itemize}
    \item Ligne \emph{Instructions so far} $\rightarrow$ métrique principale.
    \item Décomposition R / I / S / B / U / J $\rightarrow$ utile pour distinguer calcul pur vs accès mémoire (les S-type $=$ stores, les loads sont des I-type).
\end{itemize}

\paragraph{Mémoire --- lecture directe dans RARS.}
\begin{itemize}
    \item \textbf{Text Segment} : scroller jusqu'à la dernière instruction de la fonction, noter l'adresse du label d'entrée et l'adresse juste après \texttt{ret}. Différence $=$ taille du code en octets.
    \item \textbf{Data Segment} : idem pour \texttt{tab} et les constantes. On y voit le tableau avant/après le tri, utile pour vérifier.
    \item \textbf{Pile} : regarder \texttt{sp} dans la fenêtre \emph{Registers} avant l'appel, puis mettre un breakpoint sur la 1\textsuperscript{re} instruction de la fonction pour voir \texttt{sp} après le prologue. La différence $=$ octets consommés sur la pile par appel. Pour un algo itératif comme le tri à bulles, elle ne grandit jamais au-delà.
\end{itemize}

% =================================================================
\section{Résultats}
% =================================================================

\emph{Les appels des fonctions communes sont considérés négligeables. Par exemple \texttt{print\_tab} est appelée de la même manière sur tous les algorithmes.}

\begin{table}[h]
\centering
\footnotesize
\setlength{\tabcolsep}{4pt}
\begin{tabular}{lrrrrrrrrcc}
\toprule
\textbf{Algorithme} & \textbf{Instr. tot.} & \textbf{R} & \textbf{I} & \textbf{S} & \textbf{B} & \textbf{U} & \textbf{J} & \textbf{Pile (o)} & \textbf{Récursif} & \textbf{Complexité} \\
\midrule
Tri à bulles       & \textbf{2\,698} & 709 & 1\,097 & 235 & 322 & 48 & 287 & 20 & non & $O(n^2)$ \\
Tri par sélection  &                 &     &        &     &     &    &     &    &     & $O(n^2)$ \\
Tri par insertion  &                 &     &        &     &     &    &     &    &     & $O(n^2)$ \\
Tri rapide         &                 &     &        &     &     &    &     &    & oui & $O(n \log n)$ moy. \\
Tri fusion         &                 &     &        &     &     &    &     &    & oui & $O(n \log n)$ \\
Tri par tas        & \textbf{2\,481} & 925 &    830 & 144 & 385 & 48 & 149 & 16 & non & $O(n \log n)$ \\
\bottomrule
\end{tabular}
\caption{Comparaison des algorithmes de tri sur un tableau de 20 entiers.}
\label{tab:resultats}
\end{table}

% =================================================================
\section{VRAC à travailler}
% =================================================================

\subsection{Temps ($=$ nombre d'instructions)}

Comparaison directe : la colonne \emph{Total instr.} suffit. Plus petit $=$ plus rapide.

Trois lectures utiles :

\begin{table}[h]
\centering
\begin{tabular}{p{4cm} p{5cm} p{6cm}}
\toprule
\textbf{Type de comparaison} & \textbf{Formule} & \textbf{Ce que ça dit} \\
\midrule
Ratio entre algos      & $\text{total}_A / \text{total}_B$ & « A est $2{,}3\times$ plus lent que B » \\
Intensité mémoire      & $(I+S) / \text{total}$            & part des instructions qui touchent la RAM \\
Vérification théorique & $\text{total} / n^2$ ou $\text{total} / (n \log_2 n)$ & le coefficient doit rester à peu près constant si l'on refait le test avec un autre $n$ \\
\bottomrule
\end{tabular}
\end{table}

La dernière ligne est la plus intéressante pour un rapport d'archi : elle valide (ou non) la complexité théorique. Pour la bulle avec $n=20$, le coût « pur tri » $\approx 2698 - \text{print}$ instructions, divisé par $20^2 = 400$ $\rightarrow$ on obtient une constante $\sim X$. En refaisant avec $n=50$, on doit retomber sur un ordre $\sim X$.

\subsection{Mémoire}

Là, un seul chiffre ne suffit pas. Il faut distinguer deux composantes.

\paragraph{a.~Mémoire statique (code $+$ données).}

\begin{table}[h]
\centering
\begin{tabular}{p{5cm} p{9cm}}
\toprule
\textbf{Composant} & \textbf{Comment mesurer dans RARS} \\
\midrule
Taille du code de l'algo     & adresse de \texttt{ret} $-$ adresse du label d'entrée (en octets) \\
Données auxiliaires allouées & regarder \texttt{.data} du fichier : tables de travail, buffers\ldots \\
\bottomrule
\end{tabular}
\end{table}

Pour la bulle : juste du code, $0$ octet de données auxiliaires (tri in-place).
Pour le tri fusion : on allouera un tableau temporaire $\rightarrow$ $+4n$ octets de données.

\paragraph{b.~Mémoire dynamique (pile) --- c'est ça qui diffère vraiment.}

La bonne métrique n'est pas un octet absolu, c'est la fonction de $n$ :

\begin{table}[h]
\centering
\begin{tabular}{p{6cm} p{4cm} p{2.5cm} p{3cm}}
\toprule
\textbf{Algo type} & \textbf{Pile consommée} & \textbf{Pour $n=20$} & \textbf{Pour $n=1000$} \\
\midrule
Itératif (bulle, insertion, sélection) & constante (1 cadre) & $20$ o    & $20$ o \\
Récursif équilibré (fusion, tas)       & cadre $\times \log_2(n)$ & $\sim 5$ cadres & $\sim 10$ cadres \\
Récursif dégénéré (quicksort pire)     & cadre $\times n$    & $20$ cadres & $1000$ cadres ! \\
\bottomrule
\end{tabular}
\end{table}

Dans RARS on lit :
\begin{itemize}
    \item \textbf{Taille d'un cadre} $\rightarrow$ dans le prologue : \texttt{addi sp, sp, -N} $\rightarrow$ $N$ octets.
    \item \textbf{Profondeur max} $\rightarrow$ registre \texttt{sp} pendant l'exécution (breakpoint à l'intérieur de la fonction récursive), ou instrumentation \texttt{sp\_min} si besoin.
\end{itemize}

\paragraph{Résumé : ce qu'il faut présenter dans le rapport.}

Deux tableaux, pas un.

\textbf{Tableau 1 --- Temps}
\begin{center}
\begin{tabular}{llll}
Algo & Total instr. & Ratio vs plus rapide & Coefficient $\text{total} / n^2$ ou $/ (n \log n)$ \\
\end{tabular}
\end{center}

\textbf{Tableau 2 --- Mémoire}
\begin{center}
\begin{tabular}{llllll}
Algo & Code (o) & Données aux. (o) & Pile/cadre (o) & Profondeur pile & Pile max (o) \\
\end{tabular}
\end{center}

La colonne \emph{Pile max} ($=$ cadre $\times$ profondeur) est la métrique clé pour comparer l'empreinte mémoire dynamique entre algos.

\paragraph{Conclusion qualitative attendue.}

On devrait retrouver dans les chiffres :
\begin{itemize}
    \item \textbf{Bulle} : lent ($O(n^2)$) mais mémoire minimale ($20$ o fixes).
    \item \textbf{Fusion} : rapide ($O(n \log n)$) mais paye en mémoire ($4n$ o de buffer $+$ pile $\log_2 n$).
    \item \textbf{Quicksort} : rapide en moyenne, mais pile variable selon les données $\rightarrow$ risqué sur gros $n$.
\end{itemize}

\end{document}
